% Abhakseite „Das kann ich“ – Zone nur im Gesamt (\mitzone)
%% TODO Ich-kann-Sätze: Platzhalter wie in den Aufgabendateien
\begin{abhakseite}
\ifdefined\mitzone
\abhakgruppe{Kennst du schon}
\abhak{1}{Prozentwert berechnen}
\abhak{2}{Prozentsatz berechnen}
\abhak{3}{Grundwert berechnen}
\abhak{4}{Dezimalzahlen addieren und multiplizieren, Euro-Beträge auf Cent runden (zwei Dezimalen), Ergebnis mit ≈}
\abhak{5}{Dreisatz und fester Faktor („pro Monat“, „pro Tag“)}
\abhak{6}{Erhöhung um p \% als Wachstumsfaktor}
\abhak{7}{Potenz mit dem Taschenrechner}
\abhak{8}{Tabelle lesen und fortschreiben}
\abhak{9}{Prozentwert berechnen – Fehler finden}
\abhak{10}{Prozentwert berechnen – selbst rechnen}
\fi
\abhakgruppe{Einheit 1 · Jahreszins, Monats- und Tageszins}
\abhak{11}{Zinsen berechnen}
\abhak{12}{Monats- und Tageszinsen}
\abhak{13}{Überschlag: 1 \% von … , dann grob mal p}
\abhak{14}{Zinsen berechnen – Fehler finden, Begründen, Anwendung}
\abhakgruppe{Einheit 2 · Zinseszins und Guthabentabelle}
\abhak{15}{Zinseszins}
\abhak{16}{Guthaben nach n Jahren mit jährlicher Einzahlung (Tabelle statt Formel)}
\abhak{17}{Ratenkauf und Kredit mit Zinseszins (Vorrat, Verbrauchersicht)}
\abhak{18}{Zinseszins – Fehler finden, Begründen, Anwendung, Darstellungswechsel}
\end{abhakseite}
